\chapter{Antecedentes y Trabajos Relacionados} \label{cap:antecedentes}

Este capítulo presenta el contexto técnico e institucional del trabajo y la literatura sobre la que se apoya. En la Sección~\ref{sec:antecedentes} se explican los antecedentes y los conceptos fundamentales (la evolución del monitoreo hacia la observabilidad, la fatiga de alertas, los modelos de lenguaje pequeños, el ajuste fino eficiente en parámetros, la detección de anomalías en logs, las métricas de evaluación, los \textit{embeddings}, la caché semántica y el uso de modelos de lenguaje como anotadores). En la Sección~\ref{cap:trabajos_relacionados} se revisan los trabajos más cercanos y se ubica el aporte del presente estudio.

\section{Antecedentes}\label{sec:antecedentes}

\subsection{Microservicios y volumen de telemetría}\label{sec:ant_microservicios}

Las aplicaciones modernas se descomponen en servicios pequeños, desplegables de forma independiente y que se comunican mediante interfaces de red, generalmente sobre plataformas de contenedores como Kubernetes. Esta arquitectura mejora la escalabilidad y la agilidad de entrega, pero aumenta la complejidad operativa: una sola solicitud de usuario puede atravesar muchos servicios, y un fallo en uno de ellos puede propagarse al resto~\cite{wang2024_rca, uddin2026_microservice_logs}. Cada servicio emite además su propia telemetría, de modo que el volumen total crece con el número de componentes y se vuelve inviable de revisar manualmente~\cite{he2021_survey, uddin2026_microservice_logs}.

\subsubsection{El caso del corpus \texttt{ms-users}}\label{subsec:ant_corpus}

El servicio \texttt{ms-users} de la DTIC ilustra ese problema a escala institucional. Los datos disponibles para este trabajo son 309 archivos Logback/Spring Boot (\texttt{logFile.YYYY-MM-DD.HH.log}) con un total aproximado de 15.8~GB y 27{,}752{,}836 eventos. El peso del corpus se debe principalmente al registro de Spring Security en nivel \texttt{TRACE}, a los sondeos de Kubernetes a \texttt{/actuator/health} y \texttt{/actuator/prometheus} cada 5 a 15 segundos, que dejan trazas de pila de alrededor de 115 líneas, y a las redirecciones a \texttt{/login} y a los recursos estáticos.

Los eventos con valor de negocio provienen de los adaptadores del servicio hacia sistemas externos (por ejemplo, \texttt{AdUserOutAdappter}, \texttt{UsuarioAs400Service}, \texttt{UsuarioOracleService}, \texttt{GoogleUserOutAdapter} y \texttt{LdapEstudianteOutAdapter}) y registran situaciones como rechazos por política de contraseñas, usuarios no encontrados o fallos de conectividad. % TODO: confirmar con la DTIC la descripción funcional de cada adaptador y el período exacto que cubre el corpus (fechas de inicio y fin).

\subsection{Del monitoreo tradicional a la observabilidad}\label{sec:ant_observabilidad}

El monitoreo tradicional se diseñó para infraestructuras relativamente estáticas: se vigilan indicadores como CPU, memoria, disco y red, y se configuran reglas y umbrales fijos que disparan alertas. Este enfoque indica \emph{cuándo} un parámetro supera un límite, pero no explica \emph{por qué} falla un sistema, y se adapta mal a entornos con contenedores efímeros y despliegues frecuentes. % TODO: respaldar con \cite{} de fuente de observabilidad (candidatos en notes/02_observabilidad_antecedentes.md: Allam 2023; Nellutla 2026). Agregar entradas al .bib tras verificar.

La observabilidad se plantea como la capacidad de inferir el estado interno de un sistema a partir de las señales que emite, sin suponer de antemano los modos de falla. Se apoya en tres tipos de señal que deben analizarse de forma correlacionada: los logs, que son registros textuales con marca de tiempo que describen eventos, advertencias y excepciones; las métricas, que son series temporales de indicadores agregados como la latencia, la tasa de peticiones o el uso de recursos; y las trazas distribuidas, que son los recorridos causales de una petición a través de los servicios, compuestos por \textit{spans} con identificadores de traza y de span.

OpenTelemetry, promovido por la \textit{Cloud Native Computing Foundation}, se ha consolidado como estándar neutral para instrumentar, recolectar y exportar estas tres señales. % TODO: \cite{TODO} para OpenTelemetry (documentación oficial/CNCF); no hay clave en el .bib.

\subsection{Fatiga de alertas y carga operativa}\label{sec:ant_fatiga}

Los umbrales fijos en sistemas distribuidos tienden a producir un gran número de alertas ruidosas y redundantes; ante una falla en cascada, además, varios detectores disparan alertas simultáneas sobre síntomas secundarios. Esa saturación (\textit{alert fatigue}) aumenta la carga cognitiva de los equipos de operación, alarga el tiempo medio de resolución y puede llevar a ignorar fallas sutiles. Bhukar et al. señalan que la mayoría de las notificaciones que genera una plataforma de observabilidad resultan ser falsas alarmas, y que las políticas estáticas de supresión, definidas a mano, quedan obsoletas en entornos que cambian con rapidez; proponen aprender esas políticas de forma no supervisada a partir de alertas históricas y reportan una mejora de 7.39\,\% en la exactitud sobre el estado del arte en anomalías de logs y de 35.7\,\% en métricas~\cite{bhukar2024_alertsuppression}. Su trabajo actúa sobre alertas y eventos que una plataforma ya generó, no sobre líneas de log, y se validó con datos privados de la empresa, por lo que se toma aquí como respaldo del problema y no como método comparable.

En el corpus de \texttt{ms-users} el fenómeno análogo es la dominancia del ruido: el 99.2\,\% de los eventos no aporta información de negocio y, aun entre los eventos relevantes, muchos son condiciones funcionales correctamente manejadas (usuario no encontrado, política de contraseñas) y no fallas. Un mecanismo de alerta que no distinga ambas categorías reproduciría la fatiga de alertas (véase el Capítulo~\ref{cap:resultados}).

\subsection{Evolución del análisis automático de logs}\label{sec:ant_evolucion_logs}

La literatura distingue varias etapas en el análisis automático de logs~\cite{he2021_survey, jiang2021_aiops, uddin2026_microservice_logs}. La primera se basó en reglas y en \textit{parsing}: expresiones regulares escritas por expertos y analizadores que convierten los mensajes en plantillas abstrayendo las variables (por ejemplo, Drain o Spell), que requieren mantenimiento continuo y no capturan la semántica del texto. La segunda recurrió al aprendizaje automático clásico, con representaciones de frecuencia de eventos y métodos como PCA, SVM o \textit{Isolation Forest}, que ignoran en gran medida el orden y el significado de los mensajes. La tercera adoptó el aprendizaje profundo, con modelos secuenciales (LSTM o GRU, \textit{autoencoders}) como DeepLog, LogAnomaly o LogRobust, que aprenden patrones normales y marcan desviaciones. La cuarta usa modelos de lenguaje basados en \textit{Transformer}, que procesan el log como lenguaje natural y producen \textit{embeddings} semánticos más robustos ante vocabulario nuevo y cambios en el formato~\cite{murphy2024_embeddings, liu2024_logprompt, karlsen2024_benchmark}.

% TODO: DeepLog, LogAnomaly, LogRobust, Drain y Spell se mencionan en las respuestas de NotebookLM pero no tienen entrada .bib; agregar tras verificar con el notebook o con fuente primaria. Hasta entonces no se cita cada uno individualmente.

Dentro de la cuarta etapa se abre una línea que prioriza los modelos pequeños, ajustados localmente con técnicas de \textit{fine-tuning} eficiente en parámetros, como alternativa a los modelos masivos en la nube~\cite{wang2025_slm_survey, lim2025_logadreft, liu2025_lineadapter, saha2025_slm_minimal}. Esa línea es la base del presente trabajo y se desarrolla en los capítulos siguientes.

\subsection{Situación previa en la DTIC y propuesta de titulación}\label{sec:ant_dtic}

La propuesta de titulación que origina este trabajo~\cite{orellana2026_propuesta} plantea desarrollar un \textit{pipeline} de observabilidad \textit{on-premise} basado en SLM para el análisis de logs de servidores. % TODO: describir cómo se monitorean hoy los servicios de la DTIC (Grafana en el puerto 3000 aparece como servicio externo en la documentación técnica; confirmar herramientas, alertas existentes y responsables de operación).
El proyecto técnico cuenta con dos entornos de ejecución, uno con GPU NVIDIA/CUDA y otro con GPU AMD/ROCm, y con servicios locales de inferencia (Ollama y Hugging Face Transformers) y de base vectorial (Milvus Lite).

\subsection{Modelos de lenguaje pequeños}\label{sec:mt_slm}\label{cap:marco_teorico}

\subsubsection{LLM y SLM}

Los modelos de lenguaje de gran tamaño (LLM), con decenas o cientos de miles de millones de parámetros, muestran capacidades de razonamiento y comprensión contextual, pero exigen recursos de cómputo y memoria elevados para su despliegue~\cite{wang2025_slm_survey, chen2025_collab_survey}. Los modelos de lenguaje pequeños (SLM) son modelos compactos, desde millones hasta unos pocos miles de millones de parámetros (por ejemplo, DistilBERT, TinyBERT, Llama-3.2-1B/3B o Phi-3.5), que buscan equilibrar capacidad y eficiencia en entornos de recursos limitados~\cite{wang2025_slm_survey, saha2025_slm_minimal}.
% TODO: definir con precisión el umbral de parámetros que usa esta tesis para llamar "pequeño" a un modelo (el estudio incluye modelos de 0.6B a 9B parámetros; Qwen3.5-9B es el más grande).

\subsubsection{Técnicas de compresión}\label{subsec:mt_compresion}

La revisión de Wang et al.~\cite{wang2025_slm_survey} agrupa las técnicas para obtener modelos pequeños y eficientes en tres familias. La cuantización reduce la precisión numérica de pesos y activaciones (por ejemplo, de 16 bits a 8, 4 o menos bits), con ahorros de memoria de varios factores y una pérdida moderada de calidad, y se distingue la cuantización posterior al entrenamiento de la cuantización consciente del entrenamiento. La poda elimina parámetros o componentes poco relevantes, ya sean pesos individuales (no estructurada) o estructuras completas como capas, canales o cabezas de atención (estructurada). La destilación de conocimiento transfiere el comportamiento de un modelo \textit{profesor} grande a un modelo \textit{estudiante} pequeño; Yang y Pan~\cite{yang2025_lightweight} la comparan con la poda y la cuantización y concluyen que mantiene mejor la exactitud del modelo maestro, aunque su estudio no se limita a logs y sus conclusiones provienen del resumen del artículo, por lo que se citan solo como contexto general.

En este trabajo se emplean modelos ya cuantizados en formato GGUF (por ejemplo, Q4\_K\_M y Q6\_K) servidos con Ollama, junto con modelos en precisión completa cargados con Hugging Face Transformers (véase la Sección~\ref{sec:met_slm}).

\subsubsection{Ventajas del despliegue \textit{on-premise}}

Ejecutar el modelo dentro de la infraestructura institucional permite que la telemetría no salga de la organización, elimina el costo recurrente por \textit{token} de las API externas y ofrece latencias bajas adecuadas para análisis continuo~\cite{wang2025_slm_survey, pardeshi2026_smartllm, ma2026_adaptivelog}. % TODO: las cifras concretas de latencia y costo que reportan SmartLLM y AdaptiveLog están en notes/03_slm_fundamentos.md; incluirlas solo con la página/tabla exacta verificada.

\subsection{Ajuste fino eficiente en parámetros}\label{sec:mt_peft}

El ajuste fino completo de un modelo es costoso. Los métodos de ajuste fino eficiente en parámetros (PEFT) entrenan una fracción pequeña de parámetros adicionales manteniendo congelado el modelo base; la revisión de Han et al.~\cite{han2024_peft_survey} los organiza en una taxonomía que incluye los métodos reparametrizados. Nallabollu et al.~\cite{nallabollu2026_peft_cloud} evalúan adaptadores, \textit{prefix-tuning} y LoRA en entornos de bajos recursos y reportan que reducen en más de 90\,\% los parámetros entrenables conservando hasta 98\,\% de la exactitud en tareas de referencia, que no son de logs.

El método central de este trabajo es LoRA (\textit{Low-Rank Adaptation}). Hu et al.~\cite{hu2022_lora} congelan los pesos preentrenados $W$ e inyectan en cada capa del \textit{Transformer} matrices de bajo rango entrenables $A$ y $B$, de modo que la actualización es $\Delta W = BA$; parten de la hipótesis de que el cambio de pesos durante la adaptación tiene un rango intrínseco bajo. Frente al ajuste fino completo de GPT-3 175B reportan una reducción de 10\,000 veces en los parámetros entrenables y de 3 veces en la memoria de GPU, con calidad igual o mejor y sin latencia adicional en inferencia. QLoRA~\cite{dettmers2023_qlora} lleva la idea más lejos: propaga los gradientes a través de un modelo base congelado y cuantizado a 4 bits (con el tipo de dato NormalFloat de 4 bits, doble cuantización y optimizadores paginados) hacia los adaptadores LoRA, y reporta que así se ajusta un modelo de 65\,000 millones de parámetros en una sola GPU de 48~GB conservando el rendimiento del ajuste en 16 bits. Ambos artículos evalúan tareas de lenguaje general y de seguimiento de instrucciones, no la clasificación de logs, y ninguno trata el desbalance de clases.

En la literatura de análisis de logs se discuten además otras dos variantes. ReFT (\textit{Representation Fine-Tuning}) no modifica los pesos y aplica en su lugar intervenciones de bajo rango sobre las representaciones ocultas entre capas; LogADReft lo aplica a la detección de anomalías en logs~\cite{lim2025_logadreft}. LINEADAPTER inserta un adaptador convolucional 1D ligero tras la capa \textit{feed-forward} de los bloques \textit{Transformer} congelados, para capturar dependencias secuenciales locales~\cite{liu2025_lineadapter}. LogLLM, por su parte, aplica QLoRA a un Llama-3-8B para detección de anomalías~\cite{guan2024_logllm}. El \textit{fine-tuning} realizado en este trabajo utiliza LoRA con cuantización de 4 bits (véase la Sección~\ref{sec:met_finetuning}).

\subsection{Detección de anomalías en logs}\label{sec:mt_anomalias}

\subsubsection{Formulación}

Un \textit{log} es un registro textual con marca de tiempo, nivel de severidad (\texttt{TRACE}, \texttt{DEBUG}, \texttt{INFO}, \texttt{WARN}, \texttt{ERROR}), origen (\textit{logger}) y mensaje, a veces acompañado de una traza de pila. La detección de anomalías consiste en decidir, para un evento o para un conjunto de eventos, si representa un comportamiento anómalo o una falla~\cite{he2021_survey}. En este trabajo la decisión se toma sobre un \textit{batch} de eventos consecutivos: el \textit{batch} es positivo (\emph{error}) si contiene al menos un evento que constituye una falla del sistema según el criterio operativo definido en la Sección~\ref{sec:met_criterio}.

\subsubsection{Enfoques}

Los enfoques se clasifican en supervisados (requieren etiquetas), no supervisados (modelan el comportamiento normal y marcan desviaciones) y basados en modelos de lenguaje, que pueden operar en modo \textit{zero-shot} o \textit{few-shot} mediante \textit{prompting}~\cite{liu2024_logprompt, karlsen2024_benchmark} o tras ajuste fino~\cite{lim2025_logadreft, liu2025_lineadapter}. Entre los enfoques supervisados con LLM, LogLLM combina un codificador tipo BERT con un decodificador Llama-3-8B ajustado con etiquetas de anomalía~\cite{guan2024_logllm}, y Karlsen et al. ajustan decenas de modelos de clasificación de secuencias~\cite{karlsen2024_benchmark}. Los costos de usar modelos grandes en todo el flujo han motivado esquemas colaborativos, donde un SLM atiende los casos sencillos y un LLM los difíciles~\cite{ma2026_adaptivelog}, o un detector barato propone candidatos que un LLM con recuperación reevalúa~\cite{zhang2024_lograg}.

Un paso previo habitual es el \textit{parsing}, que convierte cada mensaje en una plantilla y sus variables. Sun et al. muestran que los errores de este paso se propagan de forma sistemática a la detección posterior, y que corregirlos mejora el $F_1$ de forma estadísticamente significativa, sobre todo en datos industriales~\cite{sun2026_parsing}. LogLLM, en cambio, preprocesa los mensajes con expresiones regulares sin usar un \textit{parser}~\cite{guan2024_logllm}, en la línea del enmascarado de IP, UUID y números que usa este trabajo (Sección~\ref{subsec:met_prefiltro}).

\subsubsection{Del diagnóstico a la causa raíz y la recomendación}

El análisis de causa raíz (RCA) busca identificar el componente o evento que origina una falla; las revisiones recientes cubren métodos basados en grafos, trazas y modelos de lenguaje~\cite{wang2024_rca}. Algunos trabajos combinan detección y RCA con LLM y redes bayesianas~\cite{pedroso2025_maia}. En el presente trabajo, la recomendación (causa probable y acción sugerida) se produce con un modelo de lenguaje local; no se evalúa formalmente su calidad frente a un conjunto de referencia.
% TODO: si se evalúa la calidad de las recomendaciones (p.ej. revisión experta de una muestra), agregarlo aquí y en Resultados.

\subsection{Métricas de evaluación}\label{sec:mt_metricas}

Sea $TP$, $FP$, $TN$ y $FN$ el número de verdaderos positivos, falsos positivos, verdaderos negativos y falsos negativos, tomando \emph{error} como clase positiva.

\begin{align}
\mathrm{Accuracy} &= \frac{TP + TN}{TP + FP + TN + FN}\\
\mathrm{Precision} &= \frac{TP}{TP + FP}\\
\mathrm{Recall} &= \frac{TP}{TP + FN}\\
\mathrm{Especificidad} &= \frac{TN}{TN + FP}\\
F_1 &= \frac{2\,\mathrm{Precision}\cdot\mathrm{Recall}}{\mathrm{Precision} + \mathrm{Recall}}
\end{align}

A estas se suman la exactitud balanceada, que es el promedio de la sensibilidad y la especificidad, y el coeficiente de correlación de Matthews (MCC):
\begin{equation}
\mathrm{MCC} = \frac{TP \cdot TN - FP \cdot FN}{\sqrt{(TP+FP)(TP+FN)(TN+FP)(TN+FN)}}.
\end{equation}
La curva ROC grafica la tasa de verdaderos positivos frente a la de falsos positivos al variar el umbral de decisión sobre una puntuación continua; el área bajo la curva (AUC) resume la capacidad de discriminación independientemente del umbral. En el \textit{pipeline}, la puntuación es la probabilidad $P(\text{error})$ que el modelo asigna al \textit{token} que escribe el valor booleano de su respuesta JSON, obtenida de los \textit{logits} (Hugging Face) o de los \textit{logprobs} (Ollama). 

Con clases desbalanceadas, \textit{accuracy} puede ser engañosa: un clasificador que nunca predice \emph{error} obtiene alta exactitud si los errores son raros. Chicco y Jurman muestran que la exactitud y el $F_1$ pueden dar resultados inflados con datos desbalanceados, que el $F_1$ cambia si se intercambian la clase positiva y la negativa, y que el MCC solo alcanza un valor alto si la predicción acierta en las cuatro categorías de la matriz de confusión, en proporción al tamaño de cada clase~\cite{chicco2020_mcc}. Diallo et al. llegan a una conclusión parecida al comparar el MCC, el $F_1$ y la exactitud balanceada en cuatro conjuntos de datos de salud; son conjuntos pequeños y ajenos a los logs, de modo que se citan por sus definiciones y su énfasis y no por sus cifras~\cite{diallo2025_metrics}. Respecto al AUC, Richardson et al. muestran con simulaciones que el ROC-AUC se mantiene igual al cambiar la razón entre positivos y negativos (de 1:1 a 1:99) mientras la curva de precisión-exhaustividad cambia de forma drástica para el mismo clasificador, y que el ROC-AUC solo se infla si el desbalance altera la distribución de las puntuaciones~\cite{richardson2024_roc}. De ahí que en este trabajo la precisión dependa de la prevalencia de errores y el AUC no, y que se tome el MCC, la exactitud balanceada y la precisión como criterios principales, con el AUC como medida de discriminación.

\subsection{\textit{Embeddings} y búsqueda vectorial}\label{sec:mt_embeddings}

Un modelo de \textit{embeddings} transforma un texto en un vector denso en el que textos semánticamente similares quedan próximos según la similitud coseno. Aplicado a logs, permite comparar mensajes por significado y no por coincidencia literal, y es más robusto ante vocabulario no visto y cambios de formato que los analizadores sintácticos~\cite{murphy2024_embeddings}. El proyecto utiliza Qwen3-Embedding-0.6B, de la serie Qwen3 Embedding~\cite{zhang2025_qwen3emb}, que tiene 28 capas, admite secuencias de hasta 32\,000 \textit{tokens} y produce vectores de 1\,024 dimensiones a partir del estado oculto de la última capa en el \textit{token} de fin de secuencia; sus autores declaran soporte para dimensiones personalizadas y para instrucciones, y reportan resultados de estado del arte en MTEB, un resultado de los propios autores sobre pruebas de recuperación y no sobre logs. Esa dimensionalidad puede reducirse gracias a las representaciones \textit{Matryoshka}, que entrenan un mismo vector de modo que sus primeras dimensiones ya sean un \textit{embedding} utilizable; Kusupati et al. reportan hasta 14 veces menos tamaño de \textit{embedding} en clasificación de ImageNet-1K con la misma exactitud y hasta 14 veces de aceleración en recuperación a gran escala~\cite{kusupati2022_matryoshka}, con experimentos principalmente de imágenes. El truncado solo es válido en modelos entrenados de esa manera, y requiere volver a normalizar el vector. En este trabajo los vectores se truncan a 512 dimensiones, se vuelven a normalizar y se almacenan en una base vectorial Milvus Lite; que el truncado conserve la calidad con logs enmascarados no se ha medido de forma independiente. % TODO: la nota del compañero en Qwen3 pide papers que defiendan que 512 dimensiones es un nivel apropiado en despliegues empresariales; hueco abierto.

Los \textit{embeddings} se usan de dos maneras: como clave de la caché semántica, que se describe a continuación, y para la evaluación semántica, en la que cada lote del SLM se compara con sus vecinos más próximos del \textit{ground truth} ($k$-NN), lo que permite evaluar corridas con tamaños de lote distintos al de la referencia.

\subsection{Caché semántica y reutilización de respuestas}\label{sec:mt_cache}

Llamar a un modelo de lenguaje es más caro y lento que consultar una memoria, y en muchos servicios las entradas se repiten o se parecen mucho. Una caché semántica aprovecha esto: guarda cada consulta como un \textit{embedding} junto con su respuesta y, si una consulta nueva tiene suficiente similitud con una guardada, devuelve la respuesta almacenada sin llamar al modelo. GPTCache, el sistema de código abierto que difundió la idea, se compone de un generador de \textit{embeddings}, un almacén vectorial y un evaluador de similitud, y reporta que, con acierto de caché, la respuesta llega de 2 a 10 veces más rápido que consultar el servicio externo~\cite{fu2023_gptcache}; el artículo lo escriben los desarrolladores de la herramienta, con experimentos pequeños, y reconoce que obtener buenos \textit{embeddings} para este fin es un problema abierto.

El umbral de similitud es la decisión central y tiene un compromiso inevitable. Ak et al., con 1\,000 pares de pregunta y respuesta tomados del conjunto de GPTCache y cinco modelos de \textit{embeddings}, barren el umbral de 0 a 1 en pasos de 0.01 y observan que umbrales altos dan mayor exactitud pero menos aciertos de caché, y que umbrales bajos hacen lo contrario~\cite{ak2025_siu}. Su cuerpo está en turco y sus resultados se presentan en figuras; son preguntas de chat y la exactitud se mide con un LLM como juez. Stogiannidis et al. proponen una variante que no solo devuelve respuestas guardadas: entrenan un modelo local barato con las respuestas del LLM almacenadas y deciden en cada caso entre confiar en él o llamar al LLM mediante dos criterios, la entropía de su distribución de salida y la distancia del vector de la entrada a los vecinos más parecidos de la caché; con GPT-4 como profesor, redujeron un tercio las llamadas con una caída de 0.37 puntos porcentuales en reconocimiento de intención~\cite{stogiannidis2023_cache}, sobre conjuntos de prueba balanceados que reconocen como poco realistas.

La reutilización de respuestas tiene también un efecto sobre la reproducibilidad: una caché devuelve siempre la misma respuesta guardada, mientras que recalcularla con un modelo no la garantiza. Atil et al. midieron cinco LLM configurados para ser deterministas (temperatura 0, \textit{top-p} 1 y semilla fija) en ocho tareas de opción múltiple y diez ejecuciones, y hallaron variaciones de exactitud de hasta 15\,\% entre ejecuciones; ningún modelo dio resultados repetibles en todas las tareas~\cite{atil2025_nondeterminism}. Distinguen la salida cruda de la respuesta extraída: el texto no es estable, pero la respuesta parseada lo es bastante más, y un sistema que deba interpretar la salida puede verse muy afectado cuando cambia el formato o el patrón. Atribuyen la causa, sin demostrarlo, a optimizaciones de servicio como el procesamiento por lotes continuo o la caché de prefijos; con Llama-3-8B ejecutado en GPU local, sin esas optimizaciones, obtuvieron resultados deterministas. Sus modelos son sobre todo servicios en la nube y sus tareas no son de logs.

Dentro del análisis de logs, el antecedente más cercano es LILAC, que guarda y refina en una caché las plantillas que genera el LLM al hacer \textit{parsing}, de modo que los mensajes ya vistos se resuelven sin consultarlo~\cite{jiang2024_lilac}. Los logs son muy repetitivos: en Loghub-2.0 hay más de 50 millones de mensajes y menos de 3\,500 plantillas, y LILAC reduce las consultas al LLM en varios órdenes de magnitud. La tarea es extraer plantillas y no clasificar errores, y el acierto se basa en la coincidencia de plantillas y no en la similitud de \textit{embeddings}. AdaptiveLog subraya precisamente esa limitación al señalar que este tipo de caché solo sirve para \textit{parsing} y no generaliza a otras tareas~\cite{ma2026_adaptivelog}.

\subsection{LLM como anotadores y acuerdo entre anotadores}\label{sec:mt_anotacion}

Cuando no existen etiquetas, una práctica es el \textit{pseudo-etiquetado}: varios modelos votan y se toma la mayoría. Para medir cuánto coinciden se usa el $\kappa$ de Fleiss, que extiende el coeficiente de acuerdo corregido por azar al caso de más de dos jueces~\cite{fleiss1971_kappa}:
\begin{equation}
\kappa = \frac{\bar{P} - \bar{P}_e}{1 - \bar{P}_e},
\end{equation}
donde $\bar{P}$ es el acuerdo observado medio y $\bar{P}_e$ el esperado por azar. Para interpretarlo se emplea la escala de Landis y Koch, que califica de acuerdo sustancial los valores de 0.61 a 0.80 y de casi perfecto los de 0.81 a 1.00, y que sus propios autores reconocen como arbitraria aunque útil como referencia~\cite{landis1977_kappa}.
Un $\kappa$ alto entre modelos no implica que sus etiquetas sean correctas respecto al criterio del dominio; el Capítulo~\ref{cap:resultados} documenta un caso en el que ocurre lo contrario. Por ello se complementa con revisión de expertos y se reportan intervalos de confianza de Wilson para proporciones.

La literatura respalda esa cautela. Tseng et al. evalúan LLM como anotadores en tres dominios que exigen conocimiento experto (finanzas, derecho y biomedicina), con 1\,000 instancias de cinco conjuntos etiquetados por expertos, y encuentran que técnicas de inferencia como \textit{chain-of-thought} y autoconsistencia no mejoran el desempeño, que los modelos de razonamiento superan solo marginalmente a los demás, y que una discusión entre varios LLM incrementa la exactitud y el acuerdo entre anotadores~\cite{tseng2025_expert_annotators}. Su estudio no usa logs ni modelos pequeños, así que no permite concluir cómo configurar los SLM de este trabajo. Para la construcción de etiquetas sobre logs, Sun et al. etiquetaron los conjuntos sin anotaciones con cuatro expertos en dos equipos independientes, validación cruzada y arbitraje final, y reportan $\kappa$ de Cohen entre 0.84 y 0.90~\cite{sun2026_parsing}; es el tipo de medición de acuerdo entre revisores que esta tesis no logró realizar (Sección~\ref{sec:limitaciones}).

\subsection{Notificación, recomendación e inyección de errores}\label{sec:mt_reaccion}

La reducción de la fatiga de alertas pasa por notificar solo los eventos relevantes y con contexto suficiente para actuar. La recomendación asistida por LLM complementa la alerta con una causa probable y una acción sugerida~\cite{pedroso2025_maia}. El capítulo de automatización del libro de \textit{Site Reliability Engineering} advierte que la automatización es un multiplicador de fuerza y no una panacea, y que aplicarla sin criterio puede crear tantos problemas como los que resuelve~\cite{murphy2016_sre_automation}; por eso el recomendador de este trabajo solo sugiere, y no ejecuta acciones. Para evaluar textos generados a partir de logs sin resúmenes de referencia, REFLEX propone usar un LLM como juez que califica relevancia, informatividad y coherencia~\cite{mudgal2025_reflex}; evalúa resúmenes y no recomendaciones accionables, y exigiría validar al juez contra revisores humanos, por lo que se considera solo un candidato. % TODO: \cite{} para recomendación de remediación asistida por LLM (candidato en notes/03_onpremise_chaos.md: Wang et al. 2025, Preprints.org; agregar al .bib tras verificar).

La inyección de fallos y la ingeniería del caos introducen perturbaciones controladas para evaluar la resiliencia del sistema y la capacidad de los modelos de detectarlas. MAIA, por ejemplo, usa Chaos Mesh para inyectar fallos en Kubernetes y medir \textit{recall} y $F_1$ del detector por tipo de fallo, e identifica tipos con menor desempeño~\cite{pedroso2025_maia}. En este trabajo se sigue una idea análoga a escala del \textit{batch}: se inyectan eventos de error simulados en el flujo de logs y se verifica si el \textit{batch} que los contiene se clasifica como error (véase la Sección~\ref{sec:met_inyeccion}).

\subsection{Síntesis}\label{sec:ant_sintesis}

El volumen y el ruido de la telemetría, la insuficiencia de las reglas y umbrales fijos, la fatiga de alertas, y la disponibilidad de modelos de lenguaje compactos que pueden ejecutarse localmente, ajustarse con pocos recursos y combinarse con una caché semántica definen el espacio del problema. Las secciones siguientes revisan lo que la literatura ha hecho en este espacio (Sección~\ref{cap:trabajos_relacionados}), y los capítulos posteriores describen cómo se construyó y evaluó el \textit{pipeline} para la DTIC (Capítulos~\ref{cap:metodologia} y~\ref{cap:resultados}).

\section{Trabajos relacionados}\label{cap:trabajos_relacionados}

Esta sección revisa trabajos que aplican modelos de lenguaje al análisis de logs y a la observabilidad. Se agrupan por enfoque y se cierra con una comparación y la posición del presente trabajo. Las cifras provienen de los artículos y no son comparables directamente entre sí ni con las de esta tesis, pues usan conjuntos de datos y protocolos distintos. % TODO: verificar contra el PDF original las cifras de LogPrompt, LogX, Red Flags, LogRAG, LogADReft, LINEADAPTER, LogLLM, HiPELog, MAIA, TRIDENT y AetherLog, que provienen del resumen del notebook; las de caché, métricas, LoRA y QLoRA se leyeron del PDF.

\subsection{Revisiones sistemáticas}\label{sec:tr_revisiones}

He et al.~\cite{he2021_survey} y Jiang et al.~\cite{jiang2021_aiops} sintetizan el análisis automático de logs (\textit{parsing}, detección de anomalías, diagnóstico de fallas) y su papel en AIOps. Uddin et al.~\cite{uddin2026_microservice_logs} revisan específicamente el análisis de logs de microservicios con inteligencia artificial, y Wang y Qi~\cite{wang2024_rca} el análisis de causa raíz en (micro)servicios. Wang et al.~\cite{wang2025_slm_survey} y Chen et al.~\cite{chen2025_collab_survey} cubren los SLM y los mecanismos de colaboración entre modelos grandes y pequeños, respectivamente. Estas revisiones fundamentan los conceptos de la Sección~\ref{sec:antecedentes}.

\subsection{LLM con \textit{prompting} para análisis de logs}\label{sec:tr_prompting}

LogPrompt~\cite{liu2024_logprompt} propone estrategias de \textit{prompting} (\textit{self-prompt}, \textit{chain-of-thought} e \textit{in-context learning}) para analizar logs en línea sin entrenamiento, evaluado sobre ChatGPT y sobre un modelo abierto local (Vicuna-13B). Reporta un $F_1$ promedio de 0.797 en \textit{parsing} y valores de 0.384 (BGL) y 0.450 (Spirit) en detección de anomalías, además de explicaciones valoradas por expertos. Sus límites son la baja precisión absoluta en detección y la dependencia de la API en la nube para el mejor rendimiento.

Karlsen et al.~\cite{karlsen2024_benchmark} ajustan de forma supervisada 60 modelos (BERT, RoBERTa, DistilRoBERTa, GPT-2 y GPT-Neo) para clasificar líneas de log en seis conjuntos públicos etiquetados (Apache Access, CSIC 2010, PKDD 2007, BGL, Thunderbird y Spirit). El mejor, DistilRoBERTa, alcanza un $F_1$ promedio de 0.998 y clasifica unos 135 logs por segundo. Son codificadores supervisados que clasifican líneas con las etiquetas del propio \textit{benchmark}; no son modelos generativos pequeños ni evalúan lotes, y con un $F_1$ cercano a 1 sobre conjuntos públicos su resultado contrasta con los obtenidos aquí sobre un corpus real con 2.9\,\% de lotes con error (Capítulo~\ref{cap:resultados}).

Zhang et al. (LogX)~\cite{zhang2024_logx} atacan la adaptación en línea sin reentrenar: indican al LLM pares entrada-etiqueta para corregir errores nuevos, con ChatGPT y con modelos abiertos locales (Llama3.1-8B, Qwen1.5). Reportan un $F_1$ promedio cercano a 0.80 sobre cinco conjuntos públicos (HDFS 0.842, Hadoop 0.840, OpenStack 0.747, BGL 0.830 y Thunderbird 0.813) sin datos del propio dominio, y señalan que un contexto demasiado largo en el \textit{prompt} degrada el desempeño. Anastasiadis et al. (Red Flags)~\cite{anastasiadis2026_redflags} evalúan tres LLM abiertos con cuatro estrategias de \textit{prompting} sobre una muestra balanceada de 10\,000 logs de BGL; la mejor configuración (LLaMA3.2-3B con pocos ejemplos) obtiene \textit{precision} de 89\,\%, \textit{recall} de 97\,\% y $F_1$ de 93\,\%, con 11\,\% de falsos positivos. Una muestra balanceada no refleja una prevalencia baja de errores como la de este trabajo, donde la \textit{precision} es el límite principal.

Masri et al.~\cite{masri2026_slm_severity} es el trabajo más cercano en datos y en tipo de modelo: evalúan nueve SLM y modelos pequeños de razonamiento (de 0.6B a 4B) sobre 50\,000 logs muestreados de 7.3 millones de entradas reales de \texttt{journalctl} de servidores Linux, con \textit{prompting} \textit{zero-shot}, \textit{few-shot} y recuperación aumentada (RAG, con FAISS). Qwen3-4B con RAG alcanza 95.6\,\% de exactitud, Gemma3-1B pasa de 20.25\,\% a 85.28\,\% y varios modelos de razonamiento se degradan con RAG (Phi-4-mini-reasoning supera los 228 s por log). La tarea es clasificar la severidad del log, una etiqueta de metadatos que los propios autores consideran ruidosa, y no detectar fallos por lote, y las pruebas se hicieron en lotes estáticos y no en un flujo continuo.

\subsection{Ajuste fino eficiente en parámetros para logs}\label{sec:tr_peft}

Lim et al.~\cite{lim2025_logadreft} comparan ReFT y LoRA sobre RoBERTa, GPT-2 y Llama-3 (8B) con los conjuntos HDFS, BGL, Spirit y Thunderbird. Llama3-ReFT alcanza $F_1$ de 0.9690 (BGL), 0.9950 (HDFS), 0.9192 (Spirit) y 0.9149 (Thunderbird), y supera a LoRA en 9 de 12 comparaciones; en generalización cruzada ReFT conserva un $F_1$ cercano a 0.50 mientras LoRA cae a valores cercanos a 0. El entrenamiento de ReFT es más lento que el de LoRA a rangos altos.

LINEADAPTER~\cite{liu2025_lineadapter} inserta un adaptador convolucional 1D tras la capa \textit{feed-forward} de bloques congelados (BERT, Llama-3.2-1B, Bloom-3B) y reporta $F_1$ de 0.988 (BGL), 0.980 (HDFS) y 0.953 (Hadoop) en detección, y similitud semántica de 0.959 a 0.977 con la causa raíz real, con unas cuatro veces menos parámetros y FLOPs entrenables que LoRA. Su evaluación se limita a un solo nodo de cómputo.

LogLLM~\cite{guan2024_logllm} combina un codificador BERT, que comprime cada mensaje en un vector, con un Llama-3-8B ajustado con QLoRA en tres etapas, y alcanza un $F_1$ promedio de 0.959 en HDFS (0.997), BGL (0.916), Liberty (0.958) y Thunderbird (0.966), con 150 ms por muestra; los propios autores señalan que los \textit{benchmarks} públicos no reflejan la evolución de plantillas de la industria. HiPELog~\cite{devi2026_hipelog} ajusta con LoRA un \textit{Transformer} jerárquico de dos niveles y reporta $F_1$ de 99.50\,\% en HDFS y 99.58\,\% en BGL entrenando 8\,192 parámetros, probado en lotes sin flujo continuo.

Dos trabajos generales de bajos recursos completan el panorama, sin tratar logs: Yang y Pan~\cite{yang2025_lightweight} estudian la destilación para modelos de lenguaje en entornos restringidos, y Nallabollu et al.~\cite{nallabollu2026_peft_cloud} comparan técnicas PEFT para datos en la nube. Sus resultados provienen de tareas de referencia y no de logs, por lo que solo se usan como contexto de la Sección~\ref{sec:mt_peft}.

Estos trabajos usan conjuntos públicos de benchmark (HDFS, BGL, Spirit, Thunderbird, Liberty) con etiquetas ya disponibles; el presente trabajo, en cambio, parte de un corpus real sin etiquetas y dedica un esfuerzo central a construirlas.

\subsection{Colaboración entre SLM y LLM}\label{sec:tr_colaboracion}

AdaptiveLog~\cite{ma2026_adaptivelog} combina un SLM local con un LLM: el SLM resuelve las muestras de baja incertidumbre y deriva las difíciles (alrededor del 27\,\%) al LLM, con lo que reduce los costos de API en un 73\,\% respecto de usar siempre el LLM. Reporta $F_1$ de 98.91\,\% en Thunderbird y 94.99\,\% en BGL. Mantiene la dependencia de un LLM externo o de servidores locales de gran capacidad.

LogRAG~\cite{zhang2024_lograg} propone una cascada en dos etapas: un clasificador \textit{Deep One-Class} sobre \textit{tokens} marca candidatos y un LLM (ChatGPT o Mistral-7B) los reevalúa con plantillas normales históricas recuperadas de una base vectorial. Reporta $F_1$ de 0.888 con ChatGPT y de 0.898 con Mistral en BGL, y de 0.927 en Spirit, frente a PLELog y DeepLog. Es el antecedente más cercano a la idea de reducir las llamadas al modelo con un filtro barato y una base vectorial, que en esta tesis se implementa con la caché semántica; difiere en que su detector previo se entrena con datos del sistema y requiere reentrenamiento periódico, y se evalúa sobre \textit{benchmarks} públicos.

\subsection{Reutilización de respuestas y reducción de llamadas al modelo}\label{sec:tr_cache}

Los trabajos anteriores reducen el costo con un modelo pequeño que deriva al grande los casos inciertos. Una vía complementaria es no volver a llamar al modelo cuando el caso ya se resolvió, y la literatura revisada la trata en tres contextos. En LLM de propósito general, GPTCache~\cite{fu2023_gptcache} y el estudio de Ak et al.~\cite{ak2025_siu} describen y evalúan la caché semántica y su umbral de similitud con preguntas y respuestas de chat; Stogiannidis et al.~\cite{stogiannidis2023_cache} añaden un modelo local entrenado con las respuestas guardadas y criterios de confianza, en tareas de intención y sentimiento. En análisis de logs, LILAC~\cite{jiang2024_lilac} es el único antecedente de caché encontrado, y su alcance es el \textit{parsing}.

Entre los trabajos revisados no se halló una caché semántica aplicada a decidir si un grupo de logs contiene un error. Ese es el punto donde se ubica el presente trabajo, que combina la caché por lote, la caché por evento (exacta y por similitud) y una vía rápida, y mide su efecto sobre la detección de incidentes inyectados y sobre la latencia (Capítulo~\ref{cap:resultados}). La ausencia de un trabajo previo comparable debe leerse como un resultado de esta revisión y no como una prueba de que no exista.

\subsection{Observabilidad, RCA y detección con LLM en microservicios}\label{sec:tr_observabilidad}

MAIA~\cite{pedroso2025_maia} combina un Llama-3 8B ajustado de forma no supervisada, similitud coseno entre bloques de logs y una red bayesiana para RCA, y se evalúa en Kubernetes (Sock-Shop) con inyección de fallos mediante Chaos Mesh; los $F_1$ por tipo de fallo van de 0.83 a 0.997, con menor desempeño en fallos de E/S de base de datos no relacional. TRIDENT~\cite{dutta2025_trident} integra fusión multimodal de telemetría y ranking de causa raíz con atención sobre grafos; reporta reducciones de 25\,\% en MTTR y de 30\,\% en falsos positivos diarios, pero requiere instrumentación completa de trazas. AetherLog~\cite{cui2025_aetherlog} integra un LLM con un grafo de conocimiento para el análisis de causa raíz a partir de logs, y reporta $F_1$ de 0.93 en datos de Alibaba (2\,671 fallas) y de 0.97 en China Mobile (93 fallas), con 4.83 a 4.91 s por caso; sus autores declaran como limitación que la recuperación no recorre relaciones de varios saltos. El recomendador de este trabajo no usa grafos de conocimiento ni se evalúa aún con una referencia, por lo que AetherLog se cita como contexto y no como comparación.
% TODO: otros trabajos del notebook sin entrada .bib (Wang et al. 2025 Preprints.org; TraceAI; Zhou et al. 2025; Mhatre y Shaikh 2024; Nellutla 2026; Awodele et al. 2026; LEMON). Agregar al .bib tras verificar si se desean citar.

\subsection{Modelos pequeños y despliegue local}\label{sec:tr_local}

Saha et al.~\cite{saha2025_slm_minimal} estudian el análisis de logs con SLM en entornos de recursos mínimos. SmartLLM~\cite{pardeshi2026_smartllm} propone un sistema \textit{on-premise} de modelos compactos con gobernanza y RAG, con latencia de 2 a 5 segundos y costo de API nulo, pero es una plataforma genérica sin especialización en observabilidad. Murphy et al.~\cite{murphy2024_embeddings} estudian \textit{embeddings} automáticos de mensajes de log, base del enfoque semántico usado aquí. % TODO: ampliar la descripción de Saha et al. con resultados concretos (modelos y métricas).

\subsection{Anotación con LLM y construcción de \textit{ground truth}}\label{sec:tr_anotacion}

Tseng et al.~\cite{tseng2025_expert_annotators} evalúan si LLM de gran tamaño pueden sustituir a anotadores humanos expertos en finanzas, derecho y biomedicina (1\,000 instancias de cinco conjuntos etiquetados por expertos). Las técnicas de inferencia no mejoran el desempeño (con \textit{chain-of-thought}, Claude 3 Opus baja de 64.2\,\% a 62.6\,\% de exactitud promedio y Gemini 1.5 Pro de 65.3\,\% a 63.6\,\%), los modelos de razonamiento superan solo marginalmente a los demás, y una discusión multiagente mejora la exactitud y el acuerdo entre anotadores. Los autores declaran que su evaluación es \textit{zero-shot}, limitada a tareas de clasificación con etiquetas fijas. Respalda la decisión de esta tesis de no confiar solo en el voto de modelos, aunque no trata logs ni modelos pequeños y no permite deducir cómo configurar los SLM (por ejemplo, si conviene activar el razonamiento).

Sun et al.~\cite{sun2026_parsing} estudian 16 conjuntos públicos de LogHub y cuatro conjuntos industriales propietarios, con seis \textit{parsers} (de reglas y basados en LLM) y cinco detectores. Su protocolo de etiquetado para los datos sin anotaciones usa cuatro expertos en dos equipos independientes, validación cruzada y arbitraje, con $\kappa$ de Cohen de 0.84 a 0.90. Corregir los errores de \textit{parsing} mejora el $F_1$ de forma significativa (+3.61\,\% en promedio global y +7.0\,\% relativo en los industriales), y la aumentación con logs públicos ayuda de forma dependiente del modelo y puede producir transferencia negativa; además, once ingenieros entrevistados confirman que los logs industriales son más heterogéneos que los públicos. Este trabajo comparte con ellos el uso de logs reales y de revisión experta, pero no mide el acuerdo entre revisores, y enmascara en lugar de parsear (Sección~\ref{subsec:met_prefiltro}).

Mudgal~\cite{mudgal2025_reflex} propone REFLEX, una métrica sin referencias basada en un LLM como juez para evaluar resúmenes de logs, validada con una muestra cualitativa pequeña (50 entradas) y con las limitaciones de costo y de sesgo de los \textit{embeddings} que declara el autor. Evalúa resúmenes, no recomendaciones, así que se considera solo una opción para evaluar el recomendador de esta tesis (Sección~\ref{sec:mt_reaccion}).

\subsection{Comparación y posición del presente trabajo}\label{sec:tr_comparacion}

\begin{table}[H]
\centering
\caption{Comparación cualitativa con trabajos relacionados. Los símbolos son una síntesis del autor a partir de lo descrito en las secciones anteriores.}
\label{tab:tr_comparacion}
\small
\begin{tabular}{@{}p{3.2cm}p{2.2cm}p{2.4cm}p{3.4cm}p{3.2cm}@{}}
\toprule
\textbf{Trabajo} & \textbf{Solo modelos locales} & \textbf{Datos} & \textbf{Método} & \textbf{Alerta y recomendación} \\
\midrule
LogPrompt~\cite{liu2024_logprompt} & No (API en la nube) & Benchmarks públicos & \textit{Prompting} & Explicaciones \\
AdaptiveLog~\cite{ma2026_adaptivelog} & No (híbrido) & Benchmarks públicos & SLM + LLM por incertidumbre & No \\
LogADReft~\cite{lim2025_logadreft} & Sí (ajuste local) & Benchmarks públicos & ReFT / LoRA & No \\
LINEADAPTER~\cite{liu2025_lineadapter} & Sí (ajuste local) & Benchmarks públicos & Adaptador Conv1D & RCA \\
MAIA~\cite{pedroso2025_maia} & Sí (Llama-3 8B) & Sock-Shop (sintético) & Autorregresivo + Bayes & RCA \\
LogLLM~\cite{guan2024_logllm} & No especificado & Benchmarks públicos & BERT + Llama-3-8B (QLoRA), supervisado & No \\
Masri et al.~\cite{masri2026_slm_severity} & Sí (SLM de 0.6B a 4B) & \texttt{journalctl} real de servidores Linux & \textit{Zero-shot}, \textit{few-shot} y RAG & No \\
Red Flags~\cite{anastasiadis2026_redflags} & LLM abiertos (3B a 8B) & BGL, muestra balanceada & \textit{Prompting} con pocos ejemplos & No documentado \\
LogRAG~\cite{zhang2024_lograg} & No (ChatGPT o Mistral-7B) & Benchmarks públicos (BGL, Spirit) & Detector barato + LLM con RAG & No \\
\textbf{Esta tesis} & Sí (Ollama / HF) & Corpus real DTIC & SLM + embeddings, \textit{prompt} y LoRA & Slack, recomendador, \textit{dashboard} \\
\bottomrule
\end{tabular}
\normalsize
\end{table}

% TODO: validar la tabla (columnas "Solo modelos locales" y "Datos") contra los artículos originales; la clasificación proviene del resumen del notebook.

Frente a la literatura, el presente trabajo se diferencia en cinco puntos: (i) trabaja sobre un corpus operativo real de la universidad, sin etiquetas previas, y documenta cómo se construyó el \textit{ground truth}; (ii) compara varios SLM de distintas familias y tamaños con la misma tarea, sin depender de servicios en la nube; (iii) evalúa de forma semántica con \textit{embeddings} y permite comparar tamaños de \textit{batch}; (iv) integra alerta, recomendación e inyección de errores en un mismo \textit{pipeline} con interfaz gráfica; y (v) aplica una caché semántica a la clasificación de errores en logs y evalúa su efecto, un uso que no se encontró en los trabajos revisados. Frente a los trabajos más cercanos, Masri et al.~\cite{masri2026_slm_severity} usan SLM sobre logs reales pero clasifican severidad y no fallos por lote, y Tseng et al.~\cite{tseng2025_expert_annotators} y Sun et al.~\cite{sun2026_parsing} informan sobre la anotación con LLM y la revisión experta, que aquí se combinan para construir el \textit{ground truth} de un único servicio.
